\chapter{成语梳理-详细版}


\begin{enumerate}[leftmargin=*]

    \item \textbf{可歌可泣}（kě gē kě qì）：值得歌颂，使人感动得流泪，指悲壮的事迹使人非常感动。
    
    \textbf{典故：} 语出《周易·中孚》：“得敌，或鼓或罢，或泣或歌。”明代海瑞在《方孝孺临麻姑仙坛记跋》中称：“国初方列之概，无异平原复生。追念及之，可歌可泣。”方孝孺是明初忠臣，因拒绝为朱棣起草即位诏书而被诛十族，其忠烈气节令人感佩。海瑞以“可歌可泣”形容方孝孺的事迹，后世遂以此成语形容英勇悲壮的感人事迹。
    
    \textbf{造句：} 抗日战争中无数英雄的事迹可歌可泣，永远激励着我们。

    \item \textbf{鲜为人知}（xiǎn wéi rén zhī）：很少有人知道。（“鲜”是“少”的意思）
    
    \textbf{典故：} 该成语为现代成语，出自当代作家王朔的小说《玩儿的就是心跳》：“后来伴着主人度过了那段漫长的鲜为人知的冷宫生活不知洒上了多少珍妃泪。”亦有说法认为出自蒋子龙《创作笔记》。该成语形成于20世纪后期，随着当代文学的发展而广泛使用，用以形容那些不为人知的人或事物。
    
    \textbf{造句：} 这位科学家默默无闻，其贡献长期鲜为人知。

    \item \textbf{鞠躬尽瘁}（jū gōng jìn cuì）：指小心谨慎，贡献出全部精力。
    
    \textbf{典故：} 语出三国·蜀·诸葛亮《后出师表》：“臣鞠躬尽力，死而后已。”诸葛亮（181—234），字孔明，三国时期蜀汉丞相。蜀汉建兴六年（228年），诸葛亮率军北伐曹魏，临行前上《后出师表》给后主刘禅，表达了自己鞠躬尽瘁、死而后已的决心。诸葛亮一生为蜀汉政权殚精竭虑，最终病逝于五丈原军中，真正做到了“死而后已”。此成语遂成为形容忠诚奉献精神的最高赞誉。
    
    \textbf{造句：} 周总理鞠躬尽瘁，为人民奉献了毕生精力。

    \item \textbf{当之无愧}（dāng zhī wú kuì）：承受得起某种称号或荣誉，毫不惭愧。
    
    \textbf{典故：} 出自北宋文学家欧阳修《回丁判官书》：“夫人有厚己而自如者，恃其中有所以当之而不愧也。”欧阳修（1007—1072），北宋文坛领袖，唐宋八大家之一。他在书信中阐述了一个人若内心充实、具备真正的才德，就能够坦然承受外界的赞誉而毫无愧色。明清时期该成语逐渐定型，广泛使用。
    
    \textbf{造句：} 他获得“优秀教师”称号，当之无愧。

    \item \textbf{锋芒毕露}（fēng máng bì lù）：指锐气和才干全都表现出来。多形容人气盛逞强。
    
    \textbf{典故：} 较早见于南朝·宋·范晔《后汉书·袁绍传》：“瓒示枭夷，故使锋芒挫缩，厥图不果。”现代文学作品中，华而实《汉衣冠》写道：“黄熙胤奉承地解释，想借着师友渊源、故旧情谊来笼络这位锋芒毕露的身居要位的武将。”“锋芒”本指刀剑的刃口和尖端，比喻人的锐气与才华。该成语既可作褒义（形容才华出众），也可作贬义（指人过于张扬、好表现自己），属中性成语。
    
    \textbf{造句：} 他年轻气盛，锋芒毕露，常引起同事的不满。

    \item \textbf{鹤立鸡群}（hè lì jī qún）：比喻一个人的才能或仪表在一群人里显得很突出。（褒义）
    
    \textbf{典故：} 出自东晋·戴逵《竹林七贤论》。据载，嵇绍（嵇康之子）初到洛阳时，有人对王戎说：“昨于稠人中始见嵇绍，昂昂然若野鹤之在鸡群。”“竹林七贤”是魏晋时期嵇康、阮籍、山涛、向秀、刘伶、阮咸、王戎七位名士的合称。魏晋时期政局动荡，七人常聚于竹林之下，纵酒清谈，蔑视礼法。嵇绍继承了父亲嵇康的风骨，仪表堂堂、才华出众，在人群中如鹤立鸡群，格外引人注目。
    
    \textbf{造句：} 在众多参赛选手中，他的才华鹤立鸡群。

    \item \textbf{妇孺皆知}（fù rú jiē zhī）：妇女和孩子都知道。形容广泛被人们所知晓。
    
    \textbf{典故：} 其渊源可追溯至唐代元稹《白氏长庆集序》：“王公、妾妇、牛童、马走之口无不道。”清代无垢道人《八仙全传》中亦有“果如张仙所言，形于诗歌，扮为杂剧，弄得妇孺皆知”的用法。唐代诗人白居易的诗文语言通俗易懂，流传极广，元稹在序言中描述了其作品在各类人群中广为传诵的景象，后世遂以“妇孺皆知”形容某事尽人皆知。
    
    \textbf{造句：} 《西游记》的故事在中国妇孺皆知。

    \item \textbf{家喻户晓}（jiā yù hù xiǎo）：每家每户都知道。（不能与“人人”或“知道”等连用）
    
    \textbf{典故：} 最早见于《汉书·刘辅传》：“天下不可户晓。”宋代楼钥《缴郑熙等免罪》中进一步发展为：“而遽有免罪之旨，不可以家谕（喻）户晓。”另有说法认为出自汉代刘向《列女传》中“梁节姑姊”的故事。该成语原写作“户告人晓”，意为告知每家每户，后演变为“家喻户晓”，形容名声或消息传播极广，人人皆知。
    
    \textbf{造句：} 袁隆平的事迹家喻户晓，深受人民敬仰。

    \item \textbf{锲而不舍}（qiè ér bù shě）：原意是不停地雕刻；比喻有恒心，有毅力。
    
    \textbf{典故：} 出自战国·荀况《荀子·劝学》：“锲而舍之，朽木不折；锲而不舍，金石可镂。”荀子（约前313—前238），名况，战国末期赵国人，儒家代表人物之一。《劝学》是《荀子》的首篇，系统阐述了学习的重要性和方法。荀子以雕刻为喻：如果刻几下就放弃，连朽木也刻不断；如果坚持不懈地雕刻，连金属和石头也能刻出花纹。以此强调持之以恒的重要性。
    
    \textbf{造句：} 学习需要锲而不舍的精神，才能取得成功。

    \item \textbf{目不窥园}（mù bù kuī yuán）：眼睛不看园子的景色。形容埋头读书，专心治学。
    
    \textbf{典故：} 出自《汉书·董仲舒传》。西汉大儒董仲舒“专心一意治学，三年未曾窥视园菜一眼”。董仲舒（前179—前104），西汉哲学家、经学大师。据载，他少年时酷爱学习，读书废寝忘食。其父为让他休息，在宅后修建花园，但董仲舒三年间从未入园观赏一眼，专心攻读儒家经典，终成一代儒学大师。“目不窥园”遂成为形容治学刻苦专心的经典成语。
    
    \textbf{造句：} 他目不窥园，埋头苦读，最终考上了理想的大学。

    \item \textbf{兀兀穷年}（wù wù qióng nián）：用心劳苦地一年到头这样做。
    
    \textbf{典故：} 出自唐代韩愈《进学解》：“焚膏油以继晷，恒兀兀以穷年。”韩愈（768—824），唐代文学家，唐宋八大家之首。《进学解》是韩愈任国子博士时所作的一篇论说文，以师生问答的形式阐述治学之道。“焚膏油以继晷”意为点起灯烛继续白天的学习（晷指日影），“恒兀兀以穷年”意为终年勤勉不懈。此句后来成为形容勤奋治学的经典名句。
    
    \textbf{造句：} 老教授兀兀穷年，致力于古文字研究。

    \item \textbf{呕心沥血}（ǒu xīn lì xuè）：比喻费尽心血。多用来形容事业、文艺创作等方面用心的艰苦。
    
    \textbf{典故：} 该成语由两个典故融合而成。“呕心”出自唐代李商隐《李贺小传》：李贺每日骑驴外出，背一古锦囊，遇有所得即书投囊中。其母见后叹道：“是儿要当呕出心始已耳！”“沥血”出自韩愈《归彭城》诗：“刳肝以为纸，沥血以书辞。”李贺（790—816），唐代诗人，有“诗鬼”之称。他作诗极为刻苦，常于途中觅句，日积月累，其创作方式被后人视为“呕心沥血”的典范。
    
    \textbf{造句：} 曹雪芹呕心沥血，创作了不朽名著《红楼梦》。

    \item \textbf{群蚁排衙}（qún yǐ pái yá）：指整齐地排列着。
    
    \textbf{典故：} “排衙”原指旧时官署陈设仪仗，全署属吏依次参谒长官的仪式。现代作家臧克家在《闻一多先生的说和做》中描写闻一多：“他的以十万百万字计的手稿，都是密密麻麻写得工工整整的蝇头小楷，好像群蚁排衙。”闻一多（1899—1946），现代著名诗人、学者、民主战士。他治学极为严谨，数十年如一日地从事中国古典文学研究，留下了大量工整的手稿。臧克家以“群蚁排衙”形容其手稿的整齐有序，形象地展现了闻一多一丝不苟的治学精神。
    
    \textbf{造句：} 工人们将货物码放得整整齐齐，如群蚁排衙。

    \item \textbf{鳞次栉比}（lín cì zhì bǐ）：像鱼鳞和梳子的齿一样，一个挨着一个地排列着，多形容房屋等密集。
    
    \textbf{典故：} 出自《诗经·周颂·良耜》：“穫之挃挃，积之栗栗。其崇如墉，其比如栉。”南朝宋·鲍照《咏史》诗亦有“京城十二衢，飞甍各鳞次”之句。《诗经》是中国最早的诗歌总集，收录了西周初年至春秋中期的诗歌。《良耜》是周王祭祀土神、祈求丰收的乐歌，“其比如栉”形容谷物堆积得如梳齿般整齐密集。后世以“鳞次栉比”形容建筑物排列密集整齐。
    
    \textbf{造句：} 城市里高楼大厦鳞次栉比，十分壮观。

    \item \textbf{气冲斗牛}（qì chōng dǒu niú）：形容气势之盛可以直冲云霄。（经常用来形容场面壮观）
    
    \textbf{典故：} 最早可追溯至《晋书·张华传》中“紫气冲斗牛”的典故。唐代崔融《咏宝剑》诗：“匣气冲牛斗，山形转辘轳。”宋代岳飞《题青泥赤壁》诗：“雄气堂堂贯牛斗，誓将真节报君仇。”“斗”“牛”指二十八宿中的斗宿和牛宿，代指天空。该成语原形容宝剑的光芒直射星空，后引申为形容气势磅礴、怒气冲天或精神昂扬。现代作家臧克家在《闻一多先生的说和做》中写道：“他（闻一多）说了，说的真痛快，动人心，鼓斗志，气冲斗牛，声惊天地！”
    
    \textbf{造句：} 运动场上，运动员们的斗志气冲斗牛。

    \item \textbf{微不足道}（wēi bù zú dào）：非常渺小，不值一提。（不跟意义重大的事情连用）
    
    \textbf{典故：} 出自《穀梁传·隐公二年》：“逆之道微，无足道焉尔。”《穀梁传·隐公七年》亦有类似表述。《穀梁传》是解释《春秋》的“三传”之一，相传为战国时穀梁赤所撰。书中记载了鲁国伯姬出嫁之事，因使臣地位低微，不值得特别言说，故云“无足道焉尔”。该成语反映了古代社会严格的等级观念，后泛化为形容事物渺小、不值得一提。
    
    \textbf{造句：} 这点小事微不足道，不必放在心上。

    \item \textbf{大庭广众}（dà tíng guǎng zhòng）：聚集很多人的公开场合。（不跟“公共场合”等连用）
    
    \textbf{典故：} 最早出自秦·孔鲋《孔丛子·公孙龙》：“使此人于广庭大众之中，见侮而不敢斗，王将以为臣乎？”《孔丛子》是一部记录孔子及其弟子言行的著作。书中记载了尹文与齐湣王的对话：若一个人在公众场合受人侮辱却不敢反抗，这样的“士人”是否值得任用？“广庭大众”后演变为“大庭广众”，用以指人数众多的公开场合。
    
    \textbf{造句：} 他从未在大庭广众之下发过脾气。

    \item \textbf{忍俊不禁}（rěn jùn bù jīn）：忍不住笑。（不能和“笑”连用）
    
    \textbf{典故：} 出自唐代赵璘《因话录》卷五。尚书省新做了一个存放印信的柜子，吏部郎中周戎在柜上戏作考词：“当有千有万，忍俊不禁，考上下。”赵璘是唐代小说家，《因话录》记载了唐代的逸闻轶事。“忍俊”意为含笑，“不禁”意为无法控制。周戎以幽默的方式在公物上题词，令人忍俊不禁。该成语后泛指忍不住发笑，使用时应注意不与“笑”字重复。
    
    \textbf{造句：} 小明的滑稽动作令全班同学忍俊不禁。

    \item \textbf{悲天悯人}（bēi tiān mǐn rén）：对社会的腐败和人民的疾苦感到悲愤和不平。
    
    \textbf{典故：} 出自唐代韩愈《争臣论》：“彼二圣一贤者，岂不知自安佚之为乐哉？诚畏天命而悲人穷也。”唐德宗时，谏议大夫阳城任职八年却从不进谏。韩愈作《争臣论》批评其失职，提出“君子居其位，则思死其官”的为官之道。“悲人穷”即怜悯百姓的困苦，后演变为“悲天悯人”，用以形容忧时伤世、关怀民生的情怀。
    
    \textbf{造句：} 鲁迅先生有着悲天悯人的博大胸怀。

    \item \textbf{耀武扬威}（yào wǔ yáng wēi）：炫耀武力，显示威风。（多用于贬义）
    
    \textbf{典故：} 最早可追溯至汉代士孙瑞《剑铭》：“曜威耀武。”元代乔孟符《两世姻缘》第三折：“你这般耀武扬威待怎么！”士孙瑞是东汉大臣，曾与王允共谋诛杀董卓。“耀武扬威”原指炫耀武力、显示军威，后引申为形容人得意张扬、盛气凌人的样子。《三国演义》《初刻拍案惊奇》等明清小说中均有使用。
    
    \textbf{造句：} 敌人耀武扬威地走进村庄，百姓无不愤恨。

    \item \textbf{姗姗来迟}（shān shān lái chí）：形容慢腾腾地来晚了。
    
    \textbf{典故：} 出自《汉书·外戚传·孝武李夫人》。汉武帝思念已故的李夫人，方士少翁设帐招魂，武帝遥望见一女子如李夫人之貌，作诗曰：“是耶，非耶？立而望之，偏何姗姗其来迟。”李夫人是汉武帝最宠爱的妃子，能歌善舞，容貌倾城，但早逝。汉武帝思念不已，请方士招魂。“姗姗”形容女子行走缓慢从容的样子，后“姗姗来迟”泛指来得很晚。
    
    \textbf{造句：} 会议已经开始半小时了，他才姗姗来迟。

    \item \textbf{诲人不倦}（huì rén bù juàn）：教育人极有耐心，从不厌倦。（敬辞，多形容别人）
    
    \textbf{典故：} 出自《论语·述而》：“子曰：‘默而识之，学而不厌，诲人不倦，何有于我哉？’”孔子（前551—前479），名丘，字仲尼，春秋时期鲁国人，儒家学派创始人。《论语》是记录孔子及其弟子言行的经典。“诲人不倦”体现了孔子作为教育家的伟大风范——教导他人不知疲倦，是中国教育精神的最高写照。
    
    \textbf{造句：} 张老师诲人不倦，深受学生爱戴。

    \item \textbf{不耻下问}（bù chǐ xià wèn）：不以向地位比自己低、知识比自己少的人请教为耻。（对象为不知自己的人）
    
    \textbf{典故：} 出自《论语·公冶长》：“敏而好学，不耻下问。”这是孔子对卫国大夫孔圉（谥号“文”）的评价。孔圉天资聪慧又勤奋好学，且不以向地位低于自己的人请教为耻辱，因此死后被谥为“文”。孔子以此倡导虚心求学的态度，强调学问面前人人平等。
    
    \textbf{造句：} 真正的学者都保持不耻下问的谦逊态度。

    \item \textbf{以身作则}（yǐ shēn zuò zé）：用自己的行动做出榜样。
    
    \textbf{典故：} 出自《论语·子路》：“其身正，不令而行；其身不正，虽令不从。”孔子认为，统治者如果自身行为端正，不用发布命令，百姓也会跟着去做；如果自身行为不正，即使三令五申，百姓也不会服从。这一思想后来凝练为“以身作则”，强调领导者应以身示范、率先垂范。
    
    \textbf{造句：} 班干部要以身作则，带动全班同学遵守纪律。

    \item \textbf{本末倒置}（běn mò dào zhì）：比喻把主要事物和次要事物或事物的主要方面和次要方面弄颠倒了。
    
    \textbf{典故：} 出自《礼记·大学》：“物有本末，事有终始，知所先后，则近道矣。”《大学》是儒家经典《礼记》中的一篇，系统阐述了修身、齐家、治国、平天下的思想。“本”指树根，“末”指树梢，比喻事物的根本和枝节。该成语告诫人们要分清主次，不可将轻重缓急颠倒。
    
    \textbf{造句：} 工作中不应本末倒置，要先解决主要矛盾。

\end{enumerate}